\setcounter{section}{1}
\setlength{\parindent}{0pt}  % niente rientro
\setlength{\parskip}{6pt}
\section{Requisiti non funzionali}
\subsection*{RNF-01-Compatibilità}
Il sistema deve essere compatibile e utilizzabile al pieno delle sue potenzialita
con i seguenti broswer web:
\begin{itemize}
    \item Google Crome (versione 120 o superiori)
    \item Mozilla Firefox (versione 120 o superiori)
    \item Microsoft Edge (versione 124 o superiori)
    \item Safari (versione 15 o superiori)
    \item Opera (versione 120 o superiori)
\end{itemize}
Il software deve garantire un esperienza fluida e funzionale su tutte le
piattaforme supportate indipendentemente dal broswer utilizzato, assicurando
una corretta visualizzazione delle interfacce ed il corretto
funzionamento di tutte le funzionalità presenti.
\subsection*{RNF-02-Performance}
Il sistema deve garantire tempi di risposta rapidi in modo da completare
tutte le principali operazioni (accesso, caricamento dashboard, consultazione dati)
in un tempo inferiore ai 2 secondi o nel caso di operazioni pesanti (grosse aggregazioni/grossi resoconti territoriali)
inferiore ai 5 secondi con un feedback visivo di caricamento.
Questi tempi devono essere rispettati anche durante i picchi di utilizzo dell'applicazione
che deve supportare un carico massimo di X utenti. La struttura del sistema deve 
inoltre essere scalabile in modo da essere facilemente adattata per supportare eventuali
aumenti di utenza.
\subsection*{RNF-03-Scalabilità}
L'architettura del sistema deve essere scalabile così da poter supportare una crescita
del numero di:
\begin{itemize}
    \item utenti (cittadini, enti, CER e comuni)
    \item richieste simultanee ai servizi offerti dal software
    \item dati energetici
\end{itemize}
Il sistema deve poter essere facilmente esteso a nuovi \textbf{comuni o regioni} senza modifiche
sostanziali all'infrastruttura.
\subsection*{RNF-04-Affidabilità}
Il sistema deve garantire un livello di disponibilità non inferiore al XX\%,
che corrisponde ad un downtime massimo di X ore, Y minuti, Z secondi.
Il sistema deve essere progettato per garantire accessibilità anche durante i
periodi di picco d'utilizzo. Deve inoltre essere previsto un piano di ripristino
automatico dei dati in caso di guasto (disaster recovery).
\subsection*{RNF-05-Sicurezza}
Tutti i dati devono essere memorizzati in modo sicuro nel database con tecniche di
crittografia a riposo, utilizzando algoritmi conformi agli standard AES-256
(Advanced Encryption Standard) in modalità GCM (Galois/Counter Mode)
o CBC (Cipher Block Chaining), le password devono essere salvate tramite hashing
unidirezionale con salt univoco per utente usando algoritmi moderni e
resistenti ad attacchi di forza bruta come: (da aggiungere dettagli).
\subsection*{RNF-06-Privacy e conformità normativa}
Il sistema deve essere conforme al Regolamento Europeo GDPR
(UE 2016/679) e alle normative nazionali sulla protezione dei
dei dati. Ogni trattamento deve essere documentato e devono
e tracciabile, devono essere garantiti:
\begin{itemize}
    \item il diritto all'accesso e cancellazione dei dati personali
    \item la gestione del consenso informato
    \item l'uso di accordi di trattamento dati (DPA) per soggetti esterni
\end{itemize}
\subsection*{RNF-07-Lingua}
L'interfaccia utente deve essere disponibile in lingua italiana
come predefinita, deve includere la possibilità di estensione
in altre lingue per favorire la replicabilità in territori
plurilingue o internazionali e agevolare l'utilizzo dei servizi
agli utenti di madrelingua diversa dall'italiano.
L'opzione deve essere facilmente accessibile ed applicabile in
qualsiasi momento, garantendo una traduzione automatica di tutti i contenuti
testuali e di tutte le funzionalità.
\subsection*{RNF-08-Manutentibilità}
Il codice sorgente deve essere modulare e documentato,
per consentire aggiornamenti e manutenzione evolutiva
senza compromettere la stabilità del sistema.
Devono essere previsti strumenti di versionamento (es: git)
e ambienti separati per sviluppo, test e produzione.
\subsection*{RNF-09-Usabilità}
L'interfaccia deve essere intuitiva e deve prevedere:
\begin{itemize}
    \item layout coerente e navigazione semplice
    \item icone e grafici facilmente interpretabili
    \item terminologia chiara 
    \item tempi di risposta rapidi e feedback visivo alle azioni dell'utente
\end{itemize}
così da fornire un esperienza chiara ed intuitiva anche per utenti meno
tecnici o avvezzi a queste tipi di tecnologie.
\subsection*{RNF-10-Interoperabilità}
Il sistema deve poter interagire con piattaforme esterne
(es. GSE, sistemi comunali, database open-data energetici) tramite API RESTful
standard e scambio dati in formato JSON o XML.
\subsection*{RNF-11-Tracciabilità delle operazioni}
Tutte le operazioni critiche (login, modifica dei dati, caricamento di documenti)
devono essere salvate in un log di sistema con timestamp,
ID utente e descrizione dell’azione. I log devono essere consultabili solo dagli
amministratori autorizzati e conservati secondo le politiche GDPR.
\subsection*{RNF-09-Responsività}
L’interfaccia utente deve essere completamente responsiva,
garantendo la fruibilità da smartphone, tablet e desktop.
Con obbiettivo quello di migliorare l'acessibilità e l'inclusione
dei cittadini.
\subsection*{RNF-10-Gestione errori}
Il sistema deve gestire in modo chiaro e sicuro eventuali
errori (es: connessione, inserimento dati errato, server non
disponibile) mostrando messaggi descrittivi e non
tecnici all’utente.